%%%%%%%%%%%%%%%%%%%%%%%%%%%%%%%%%%%%%%%%%%%%%%%%%%%%%%%%%%%%%%%%%%%%%%
% Lớp 4 | Chủ đề 1. Ôn tập và bổ sung | Biểu thức chứa chữ
% Dạng 4D1?4-3: Chu vi, diện tích viết bằng biểu thức chứa chữ
%%%%%%%%%%%%%%%%%%%%%%%%%%%%%%%%%%%%%%%%%%%%%%%%%%%%%%%%%%%%%%%%%%%%%%


% nguồn: 2026-11-18_lop4_chu-de-1-tu-luan (nhập 18/11/2026)
\begin{bt}%[4D1N4-3]%[ma:d3839e56]
% nguon: Toán tư duy lớp 4 – Bài 4 SGK Toán 4 KNTT
	Một hình vuông có cạnh là $a$. Viết biểu thức tính chu vi $P$ của hình vuông. Tính $P$ khi $a = 15$ cm.
	\loigiai{
		Hình vuông có $4$ cạnh bằng nhau nên $P = a \times 4$.\\
		Với $a = 15$ cm: $P = 15 \times 4 = 60$ (cm).
	}
\end{bt}

% nguồn: 2026-11-18_lop4_chu-de-1-tu-luan (nhập 18/11/2026)
\begin{bt}%[4D1H4-3]%[ma:3614d3d3]
% nguon: Toán tư duy lớp 4 – Bài 4 SGK Toán 4 KNTT
	Một hình chữ nhật có chiều dài $a$ và chiều rộng $b$ (cùng đơn vị đo). Viết biểu thức tính chu vi của hình chữ nhật. Tính chu vi và diện tích khi $a = 24$ cm, $b = 16$ cm.
	\loigiai{
		Chu vi: $P = (a + b) \times 2$.\\
		Với $a = 24$ cm, $b = 16$ cm: $P = (24 + 16) \times 2 = 80$ (cm).\\
		Diện tích: $24 \times 16 = 384$ (cm$^2$).
	}
\end{bt}

% nguồn: 2026-11-18_lop4_chu-de-1-tu-luan (nhập 18/11/2026)
\begin{bt}%[4D1V4-3]%[ma:554e66da]
% nguon: Toán tư duy lớp 4 – Bài 4 SGK Toán 4 KNTT
	Một hình chữ nhật có chiều rộng $a$ cm, chiều dài gấp $3$ lần chiều rộng. Viết biểu thức tính chu vi hình chữ nhật theo $a$. Tính chu vi và diện tích khi $a = 9$.
	\loigiai{
		Chiều dài là $a \times 3$.\\
		Chu vi: $(a + a \times 3) \times 2 = a \times 4 \times 2 = a \times 8$.\\
		Khi $a = 9$: chu vi là $9 \times 8 = 72$ (cm); chiều dài $9 \times 3 = 27$ (cm).\\
		Diện tích: $27 \times 9 = 243$ (cm$^2$).
	}
\end{bt}

% nguồn: 2026-11-18_lop4_chu-de-1-tu-luan-bo-sung (nhập 18/11/2026)
\begin{bt}%[4D1H4-3]%[ma:16574e36]
% nguon: Toán tư duy lớp 4 – Bài 4 SGK Toán 4 KNTT
	Một hình vuông có chu vi là $P$. Viết biểu thức tính độ dài cạnh hình vuông theo $P$. Tính cạnh và diện tích hình vuông khi $P = 36$ cm.
	\loigiai{
		Chu vi bằng cạnh nhân $4$, nên cạnh $= P : 4$.\\
		Khi $P = 36$ cm: cạnh là $36 : 4 = 9$ (cm).\\
		Diện tích: $9 \times 9 = 81$ (cm$^2$).
	}
\end{bt}
